\documentclass[12pt]{article}
\usepackage[utf8]{inputenc}
\usepackage[T1]{fontenc}
\usepackage{hyperref}
\title{Western Lalu Language Resources for K-12 Curriculum Development}
\author{Compiled by Claude Sonnet 5.5 (claude-sonnet-5-5) for Andr\'as Kornai}
\date{}
\begin{document}
\maketitle

\section{Introduction}
Western Lalu (ISO 639-3 ywl; Glottocode west1506) is a Loloish (Sino-Tibetan) language of western Yunnan, China. Its speakers are officially classed as Yi. It belongs to the Lalo or ``Western Yi'' cluster, which is split across several ISO codes (ywt, yik, yit and ywl). Sources disagree on boundaries and speaker numbers. Wikipedia gives about 44,000 West Lalo speakers in Yongping, Yangbi and Longyang counties. Joshua Project lists a related group, the Yangliu Lalu, at about 48,000, in the Dehong and Baoshan areas. The searches found no school materials in the language, and little material that is specific to ywl.

\section{Primary Educational Resources}
The searches found no primers, textbooks or school programmes specific to Western Lalu.

\section{Linguistic and Reference Materials}
\begin{itemize}
\item Wikipedia, Lalo language. It covers the Lalo cluster, including the Western varieties. \url{https://en.wikipedia.org/wiki/Lalo_language}
\item Max Planck EVA Channumerals database: Western Lalu, Yunnan (numeral data). \url{https://lingweb.eva.mpg.de/channumerals/Lalu-Western.htm}
\end{itemize}

\section{Religious and Community Resources}
These are outsider profiles of related Lalu groups. They are not language resources.
\begin{itemize}
\item Joshua Project, Yangliu Lalu in China: \url{https://joshuaproject.net/people_groups/18536}
\item Asia Harvest profile, Lalu Yangliu: \url{https://people-groups.asiaharvest.org/people-group-profiles/yunnan/lalu-yangliu.pdf}
\item Asia Harvest profile, Lalu Xinping: \url{https://people-groups.asiaharvest.org/people-group-profiles/yunnan/lalu-xinping.pdf}
\item Asia Harvest profile, Lalu Xuzhang: \url{https://people-groups.asiaharvest.org/China/chinaPeoples/L/LaluXuzhang.pdf}
\end{itemize}

\section{Digital and Audio Resources}
No audio, video or software for Western Lalu was found.

\section{Curriculum Development Context}
The searches found no evidence of Western Lalu being taught in schools. Schooling in the area is in Mandarin. Yi-language education in China mostly uses standardized Yi (Nuosu) in Sichuan, and nothing found links that to Western Lalu. Asia Harvest describes the Yangliu Lalu as assimilating rapidly. A curriculum would need new documentation work: a community-approved orthography, word lists and recordings, and local teacher training. Chinese-language Yi dictionaries and the Lalo fieldwork literature, such as the Wang and Zhao work cited on Wikipedia, are the likely starting points.

\end{document}